% Pursuing Stacks (cotes 134-2 to 134-8): archive cote 134-7, site page 18 (PDF page 19), typescript page 460.
% typescript. Draft compared: ps6.tex:834-896.
% AI-transcribed by Claude Opus 5.5 (claude-opus-5-5, adaptive thinking, effort high), Message Batches API, batch ps3 (msgbatch_01XdFmuVVBAcgc7EFxdtCqsj), 2026-11-12; not yet checked by eye.
\scanpage{134-7/18}
\mspage{460}

26.8.

112. I am continuing the ``wishful thinking'' abo\over{\unsure{u}}{u}t schematization of
homotopy types - a wel\deleted{l}come break in the ``overall review'' on lineariza-%
tion \xout{\unsure{of}} (in the context of the modelizer (Cat)), which had been getting
a little fastidious lately !

I'll admit, as one firm hold in all the wishfulness, that in the
``Eilenberg-Maclane case'' of p.453, when moreover $k=\underline{Z}$ and the components
$X_i$ of the semisimplicial unipotent bundle $X_*$ are projective, the two
maps (w) of page 454 are indeed quasiisomorphisms. From this should
follow the similar statement, when $X_*$ comes from a chain complex of
k-modules $M.$ with projective coefficients, by reducing to the case
when only a finite number of components of $M.$ are non zero (by suitable
\margin{\unsure{!}} passage to the limit), and then by \uline{induction} on the number of these
components, using the fact that the three terms in (w) depend on $X_*$ in
a ``multiplicative'' way, namely direct sums being transformed into tensor
products. (NB Under the assumptions of projectivity made, we may as
well express the quasi-isomorphis\unsure{ms}(w) we start with as being semisimpli-%
cial \uline{homotopisms}, and remark that comp\inserted{o}nentwise tensorproduct of such
homotopisms is again one.) From this, using the relevant spectral se-%
quences in homology, should follow that the maps (w) are still quasi-iso-%
morphisms, whenever $X_*$ can be ``unscrewed'' (``dévissé'') as a finite succes-%
sive fibering with fibers of the type $M_*$ as above. Another passage to
the limit will yield the same result for an infinite dévissage, provided
the fibers $M(i)_*$ ($i=1,2, \ldots$) are ``way-out'' i.e. for given n, only
a finite number of components $M(i)_n$ are non zero (it amounts to the
same to \unsure{demande|dehande} that the sequence $M(i).$ of corresponding chain complexes
be ``way out'' ). This will give already, it seems, a fair number of
cases when (w) are quasi-isomorphisms. (Admittedly, working this out will
involve a fair amount of work, especially for getting the relevant
properties of $L.$ and $L_{pt}$ , which should mimick very closely the known
ones for usual linearization, including spectral sequences or, more
neatly, transitivity isomorphisms in the relevant derived categories...)
The main point here is that those special types of $X_*$ 's (we may call
them \uline{Kan-Postnikov} \uline{complexes} in $U(\over{k}{k})$ ) are enough in order to modelize,
via the corresponding semisimplicial sets $X_*(\underline{Z})$ , arbitrary \xout{\unsure{hom}} pointed
homotopy types with abelian $\hand{\pi}_1$ . This is seen of course using Postnikov

\margin{18}
